\documentclass[10pt,a4paper]{amsart}
\usepackage[margin=19mm]{geometry}
\usepackage{amsmath,amssymb,mathtools}
\usepackage[scr=rsfs,cal=euler]{mathalfa}
\usepackage{xfrac}
\usepackage[unicode,hidelinks,bookmarks=false]{hyperref}
\newcommand{\ie}{i.e.\ }
\newcommand{\cf}{cf.\ }
\newcommand{\resp}{respectively\ }
\newcommand{\Subsection}{\subsection}
\providecommand{\nobreakdash}{\nobreak\hbox{-}}
\setlength{\emergencystretch}{2em}
\multlinegap0pt
\usepackage{amssymb}
\usepackage[shortlabels]{enumitem}
\usepackage{mathtools}
\newtheorem{theorem}{Theorem}
\newtheorem{proposition}{Proposition}
\newtheorem{lemma}{Lemma}
\def\Fbb{\mathbb{F}}
\def\l({\left(}
\def\lo{\left\{}
\def\ro{\right\}}
\def\lv{\left\vert}
\def\r){\right)}
\def\rv{\right\vert}
\title{Packing sets under finite groups via algebraic incidence structures}
\author{Norbert Hegyv\'ari, Le Quang Hung, Alex Iosevich, Thang Pham}
\date{Source: arXiv:2411.05377v4 (CC BY 4.0)}
\begin{document}
\maketitle

\noindent\emph{Source and licence.} Packing sets under finite groups via algebraic incidence structures. Version arXiv:2411.05377v4 (CC BY 4.0), distributed by arXiv under the Creative Commons Attribution 4.0 International licence, http://creativecommons.org/licenses/by/4.0/, as recorded in the arXiv licence metadata for this exact version and shown on the abstract page https://arxiv.org/abs/2411.05377v4. Full licence notice: \texttt{licenses/arxiv-2411.05377-CC-BY-4.0.txt}.\par\smallskip

\noindent\textit{Editorial scope.} Counted unit: the article's theorem H-inci (source0.tex lines 1284--1287). The article's complete original proof of it is delivered with it (source0.tex lines 1294--1360, one \texttt{proof} environment of 67 lines, which the article places away from the statement). No mathematics is written, repaired or re-derived: the statement and the proof are the article's own lines, in the article's own notation, and no retained line is substituted or re-flowed.  No further source-local context is needed: every cross-reference of every retained line resolves inside the two delivered spans.  Nothing was omitted from any retained span, so the package carries no omission record.  No retained line contains a citation, so the wrapper carries no bibliography and no bibliography entry.  No retained line contains a figure environment, an image-inclusion command or drawing code, so no image is delivered and the package declares no asset dependency.  Editorial formatting only, in addition: an AMS \texttt{amsart} preamble that restates the article's own declarations and macros used by the retained lines, namely the environments \texttt{theorem}, \texttt{proposition}, \texttt{lemma} on separate counters, together with the standard packages those lines need.  The retained environments print in this document as Theorem 1, Proposition 1, Lemma 1 in order of appearance, whereas the article prints the same items with the numbering of its own full document.  The complete native source of the article as distributed in the arXiv e-print for this exact version is bundled unchanged as \texttt{sources/168200/source0.tex}.
\begin{theorem}\label{H-inci}Let $\epsilon\ge 0$.
    Let $X\subset \mathbb{H}_1(\Fbb_p)$ and $P=A\times B\subset \mathbb{F}_p^6$. Assume that all points in $A$ and $B$ have the third coordinate non-zero, and for each $(y, z)\in \mathbb{F}_p^2$, we have $|\pi_{23}^{-1}(y, z)\cap B|\le p^{1-\epsilon}$, then the number of incidences between $X$ and $P$ satisfies
    \[\left\vert I(P, X)- \frac{|P||X|}{p^3} \right\vert\le p^{\frac{3-\epsilon}{2}}|P|^{\frac{1}{2}}|X|^{\frac{1}{2}}.\]
\end{theorem}

\begin{proof}[Proof of Theorem \ref{H-inci}]
By repeating the argument as in the case of $SL_2(\mathbb{F}_p)$, we have 
    \begin{align*}
        I(P,X) & = \sum_{\substack{(x,y )\in P \\ \theta \in X}} 1_{\theta x=y } =\frac{1}{p^3} \sum_{m\in \Fbb_p^3} \sum_{\substack{(x,y)\in P, \\ \theta \in X}} \chi \l( m\cdot \l( \theta y-x \r) \r) \\
        & = \frac{\lv P\rv \lv X\rv}{p^3} + \frac{1}{p^3} \sum_{ m\in \Fbb_p^3 \setminus \lo (0,0, 0)\ro } \sum_{\substack{ (x,y)\in P ,\\ \theta \in X }} \chi \l( m \cdot \l( \theta y -x \r) \r) \\
        & = \frac{\lv P\rv \lv X\rv}{p^3} + p^3\sum_{m\ne (0, 0, 0)} \sum_{\theta \in X} \widehat{P} \l( -m,\theta^t m\r).
    \end{align*}
By using the Cauchy--Schwarz inequality, one has 
\begin{align*}
    &\sum_{\theta\in X}\sum_{m\ne (0, 0, 0)}\widehat{P}(-m, \theta^tm)\le |X|^{\frac{1}{2}}\left(\sum_{\theta\in \mathbb{H}_1(\Fbb_p)}\sum_{m_1, m_2\ne (0, 0, 0)}\widehat{P}(-m_1, \theta^tm_1)\overline{\widehat{P}(-m_2, \theta^tm_2)}\right)^{\frac{1}{2}}.
\end{align*}
We now observe 
\begin{align*}
    &\sum_{\theta\in \mathbb{H}_1(\Fbb_p)}\sum_{m_1, m_2\ne (0, 0, 0)}\widehat{P}(-m_1, \theta^tm_1)\overline{\widehat{P}(-m_2, \theta^tm_2)}\\&=\frac{1}{p^{12}}\sum_{\theta}\sum_{m_1, m_2\ne (0, 0, 0)}\sum_{(x_1, y_1), (x_2, y_2)}P(x_1, y_1)P(x_2, y_2)\chi(-m_1x_1+\theta^tm_1y_1)\chi(m_2x_2-\theta^tm_2y_2)\\
    &=\sum_{m_1, m_2}-\sum_{m_1=(0, 0, 0), m_2\ne (0, 0, 0)}-\sum_{m_1\ne (0, 0, 0), m_2=(0, 0, 0)}-\sum_{m_1=m_2=(0, 0, 0)}\\
    &=I-II-III-IV.
\end{align*}
We now estimate each term separately. 

Regarding $IV$, it is clear that it is equal to $|P|^2/p^9$. The terms $II$ and $III$ are the same, and we can see that 
\[II=\frac{|P|}{p^{12}}\sum_{\theta}\sum_{m_2\ne (0, 0, 0)}\sum_{(x_2, y_2)}P(x_2, y_2)\chi(m_2(x_2-\theta y_2)),\]
which can be written as 
\[\frac{|P|(p^2-1)}{p^{12}}\sum_{\theta}\sum_{(x_2, y_2)\in P, x_2=\theta y_2} 1-\frac{|P|}{p^{12}}\sum_{\theta}\sum_{(x_2, y_2)\in P, x_2\ne\theta y_2}1.\]
So, $-II\le \frac{|P|^2}{p^9}$.
Regarding $I$,
\begin{align*}
    &I=\frac{1}{p^{12}}\sum_{\theta}\sum_{m_1, m_2\in \mathbb{F}_p^3}\sum_{(x_1, y_1), (x_2, y_2)}P(x_1, y_1)P(x_2, y_2)\chi(m_1(\theta y_1-x_1))\chi(m_2(\theta y_2-x_2))\\
    &=\frac{1}{p^6}\sum_{\theta}\sum_{(x_1, y_1), (x_2, y_2)}P(x_1, y_1)P(x_2, y_2)1_{\theta y_1=x_1}1_{\theta y_2=x_2}.\\
\end{align*}
To proceed further, we need to count the number of quadruples $(x, y, z)\in B, (x', y', z')\in A, (u, v, w)\in B, (u', v', w')\in A$ such that there exists $\theta=[a, b, c]\in \mathbb{H}_1(\Fbb_p)$ and $\theta (x, y, z)=(x', y', z')$, ~$\theta(u, v, w)=(u', v', w')$.

For such quadruples, one has 
\[yw'+zv'=y'w+vz', ~z(u'-u)+w(x-x')=a(vz-wy), ~z=z', w=w'.\]
From here, as in the case of the special linear group, we need to make use of a number of preliminary lemmas, which will be proved later.
\begin{proposition}\label{propo2.1}
Given $A, B\subset \mathbb{F}_p^3$, let $N$ be the number of quadruples $(x, y, z)\in A, (x', y', z')\in B, (u, v, w)\in A, (u', v', w')\in B$ such that 
\[yw'+zv'=y'w+vz', ~z=z', w=w'.\]
Suppose in addition that for each $(y, z)\in \mathbb{F}_p^2$, we have $\pi_{23}^{-1}(y, z)\cap A$ or $\pi_{23}^{-1}(y, z)\cap B$ are of sizes at most $p^{1-\epsilon}$, then we have 
\[N\le \frac{|A|^2|B|^2}{p}+p^{3-2\epsilon}|A||B|.\]
\end{proposition}
\begin{proposition}\label{propo2.2}
    Given $A, B\subset \mathbb{F}_p^3$, let $N'$ be the number of quadruples $(x, y, z)\in B, (x', y', z')\in A, (u, v, w)\in B, (u', v', w')\in A$ such that 
\[yw'+zv'=y'w+vz', ~vz-wy=0, ~zu'+xw'-z'u-x'w=0, ~w=w',~z=z'.\]
Then $N'\le p|A||B|$.
\end{proposition}
\begin{lemma}\label{lm23}
    Given $(x, y, z), (x', y', z'), (u, v, w), (u', v', w')$ in $\mathbb{F}_p^3$ with $w\ne 0$, $z\ne 0$, and $yw'+zv'=y'w+z'v$. Let $N''$ be the number of matrices $\theta=[a, b, c]\in \mathbb{H}_1(\Fbb_p)$ such that 
\[\theta(x, y, z)=(x', y', z')\]
and 
\[\theta(u, v, w)=(u', v', w').\]
\begin{itemize}
    \item If $vz-wy\ne 0$, then $N''=1$. 
    \item If $vz-wy=0$ and $z(u'-u)=w(x-x')$, then $N''=p$.

\item If $vz-wy=0$ and $z(u'-u)\ne w(x-x')$, then $N''=0$.
\end{itemize}
\end{lemma}


So, combining Propositions \ref{propo2.1}, \ref{propo2.2}, and Lemma \ref{lm23} gives
\[I\le \frac{|P|^2}{p^3}+p^{3-\epsilon}|P|.\]
In other words,
\[I-II-III-IV\ll \frac{1}{p^6}\cdot p^{3-\epsilon}|P|,\]
and 
 \[\left\vert I(P, X)- \frac{|P||X|}{p^3} \right\vert\le p^{\frac{3-\epsilon}{2}}|P|^{\frac{1}{2}}|X|^{\frac{1}{2}}.\]
 This completes the proof.
\end{proof}

\end{document}
